% !TEX program = xelatex
\documentclass[12pt,a4paper]{ctexart}

% === 页面布局 ===
\usepackage[top=2.54cm, bottom=2.54cm, left=3.18cm, right=3.18cm]{geometry}

% === 行距 ===
\usepackage{setspace}
\linespread{1.4}

% === 字体 ===
\setCJKmainfont{SimSun}[AutoFakeBold=2.5]
\setCJKsansfont{SimHei}
\setCJKmonofont{SimSun}
\setmonofont{Consolas}[Scale=0.82]

% === 标题格式（一级"一、"黑三居中；二级"（一）"；三级"1."）===
\ctexset{
  section = {
    name = {,、},
    number = \chinese{section},
    format = \zihao{3}\heiti\centering,
    aftername = \quad
  },
  subsection = {
    name = {（,）},
    number = \chinese{subsection},
    format = \zihao{4}\heiti,
    aftername = \quad
  },
  subsubsection = {
    name = {,.},
    number = \arabic{subsubsection},
    format = \zihao{-4}\heiti,
    aftername = \quad
  }
}

% === 必备宏包 ===
\usepackage{amsmath,amssymb,amsthm}
\usepackage{booktabs}
\usepackage{graphicx}
\usepackage{float}
\usepackage[colorlinks=true, linkcolor=black, citecolor=black, urlcolor=black]{hyperref}
\usepackage{caption}
\usepackage{enumitem}
\usepackage{fvextra}
\usepackage{array}

\captionsetup{font={small}, labelsep=quad}
\pagestyle{plain}

% 附录代码排版设置
\fvset{frame=single, numbers=left, numbersep=6pt, fontsize=\zihao{-5},
       baselinestretch=1.05, tabsize=4, breaklines=true, breakanywhere=true}

\begin{document}

% ================= 摘要页 =================
\begin{titlepage}
\centering
{\zihao{3}\heiti “互联网+”时代的出租车资源配置：\\[0.3em]基于分时段双边市场均衡的测量、评价与优化\par}
\vspace{0.5cm}
{\zihao{4}\heiti 摘\hspace{2em}要\par}
\vspace{0.2cm}

\begin{minipage}{0.95\textwidth}
\zihao{-4}
\setlength{\parindent}{2em}
针对“互联网+”时代的出租车资源配置问题，本文以北京为案例建立分时段双边市场广义成本均衡模型，将供求匹配测量、补贴评价与补贴优化统一于同一均衡引擎，三问分别作为引擎的测量、评价与优化模式运行：乘客需求随广义成本（运费加候车时间成本）弹性变化，司机供给随收入弹性变化，候车时间经供需比负反馈形成不动点均衡，反馈增益 $g=0.136$（基准情形最大值）$<1$ 保证均衡唯一。

针对问题一，构建“时间—空间”双维度指标体系：时间维度提出匹配度、打车难指数、候车时间与需求加权综合匹配指数 SDMI，按撮合与物理双口径计量；空间维度以万人拥有量功效系数与城市内三区供需比刻画。结果表明：北京打车难呈显著时段结构，早晚高峰候车 19.3/18.2 分钟、打车难指数 0.41/0.36、应答率仅 0.59/0.64，平峰与夜间则供给过剩；空间上北京万人拥有量 31.2 辆/万人，超出住建部标准 56\%，打车难的根源是时空错配而非总量不足。

针对问题二，以反事实识别分离平台信息效应与补贴资金效应。2000 次蒙特卡洛仿真表明，平台信息效应使高峰候车下降 2.23 分钟（90\% 区间 $[-3.56,-1.08]$，方向概率 100\%）；而滴滴 2014 年 1 月“乘客 10 元加司机 10 元”方案使高峰未满足需求不降反升 3.85 万单，日补贴 3458.7 万元反而使日失配成本恶化 22.5 万元。机理是需求弹性 0.62 大于供给弹性 0.3，高额双侧补贴在高峰反噬打车难（方向概率 99.65\%）。2014 年 5 月后纯司机侧低补贴方案为实际方案中最优，日失配成本下降 76.5 万元，成本效益比 0.18。

针对问题三，以日失配成本最小化为目标、日补贴预算 100 万元为约束求解八维非凸规划，22 个起点收敛一致（相对标准差 $4.9\times10^{-14}$），并经 39 万点双层网格验证（相对差距 0.21\%/0.11\%）。最优方案为纯司机侧、纯高峰补贴（早高峰 2.36 元/单、晚高峰 2.87 元/单），日失配成本降至 687.2 万元（改善 2.64\%），高峰候车下降 2.1\%，高峰打车难指数下降 6.1\%；预算 $[0,300]$ 万元全区间结构恒定，成本比值越出临界区间 $(2.0,2.5)$ 时结论仍稳健。

\vspace{0.3cm}
\noindent{\heiti 关键词：}双边市场均衡；反事实识别；补贴反噬；非凸优化；蒙特卡洛仿真
\end{minipage}
\end{titlepage}

% ================= 一、问题重述 =================
\section{问题重述}

\subsection{问题背景}

出租车是城市公共交通的重要补充。长期以来，乘客在路边扬招候车、司机沿街巡游找客，信息不对称使“乘客打不到车”与“司机空驶率高”并存，供求在时间与空间上的错配成为城市管理的难题。随着“互联网+”行动推进，滴滴、快的等打车软件平台兴起，通过移动互联网实时撮合供求双方，并在 2014 年前后掀起大规模补贴竞争。补贴是否缓解了打车难、应当如何设计补贴方案，成为兼具学术与政策价值的问题。本题以北京为主要案例，要求对出租车资源的供求匹配程度进行测量，对平台补贴方案作出评价，并设计新的补贴方案。

\subsection{需要解决的问题}

将赛题要求提炼为如下三个明确的数学任务。

\textbf{问题一（测量）}：建立合理的指标体系，量化刻画不同时代背景下（传统扬招与平台撮合）出租车资源在“时间—空间”两个维度上的供求匹配程度，回答“打车难到底难在何时、何地、难到什么程度”。

\textbf{问题二（评价）}：针对滴滴、快的等公司在 2014 年前后推出的多轮实际补贴方案，建立可识别的评价框架，分离平台本身的信息效应与补贴资金效应，定量回答“各公司的补贴是否缓解了打车难、缓解或恶化了多少”。

\textbf{问题三（设计）}：站在一个新打车软件平台的立场，在有限日补贴预算约束下，以社会失配成本最小化为目标，优化补贴在乘客侧与司机侧、在不同时段之间的分配结构，并论证方案的合理性与稳健性。

\subsection{名词解释与数据说明}

题目涉及的专有名词与本文数据口径如表 \ref{tab:terms} 所示。本题无附件数据，所用公开数据的检索来源在正文首次出现处随文标注。

\begin{table}[H]
\centering
\caption{名词解释与数据说明}
\label{tab:terms}
\zihao{-5}
\begin{tabular}{p{3.2cm}p{10.5cm}}
\toprule
名词 & 说明 \\
\midrule
打车软件平台 & 通过移动互联网实时撮合乘客叫单与司机接单的中介系统，本文以信息效率乘数 $\kappa$ 刻画其撮合效率 \\
补贴方案 & 平台向乘客端（每单立减）与司机端（每单奖励）发放的资金组合，分别记为 $s_p$、$s_d$（元/单） \\
供求匹配 & 同一时段内乘客需求与司机供给在数量上的平衡状态，从成交、候车、空驶多角度计量 \\
广义成本 & 乘客完成一次打车承担的总成本，等于货币运费加候车时间折算成本（元/单） \\
数据口径 & 无附件；官方统计、新闻报道与研究报告等公开检索数据在正文随文标注来源，构造标定数据在数据节显式声明 \\
\bottomrule
\end{tabular}
\end{table}

% ================= 二、问题分析 =================
\section{问题分析}

三个问题看似分属测量、评价与设计三类任务，实质是同一市场系统在三种输入下的三种输出：问题一是基准输入下的状态测量，问题二是给定政策输入下的反事实比较，问题三是政策作为决策变量的优化。若能建立一个同时容纳乘客需求行为、司机供给行为与平台撮合技术的统一市场均衡模型，则三问可在同一引擎上求解，指标口径、参数标定与验证标准全程一致。基于这一判断，本文先构建分时段双边市场广义成本均衡引擎，再让引擎分别以“测量模式、评价模式、优化模式”运行。以下分述各问的核心难点与方法路线。

\subsection{时空双维度供求匹配的测量问题分析}

问题一的难点有三。其一，“匹配程度”没有现成定义：供给大于需求时乘客易打车但司机空驶，供给小于需求时反之，单一比值无法同时刻画两种失衡，需要构造对两种失衡都敏感的指标族，并区分“高峰短缺”与“平峰过剩”两种性质相反的状态。其二，平台出现前后匹配程度的变化必须可比：打车软件提高的是撮合效率而非物理运力，若把撮合加快带来的成交增量全部记为“资源变多”，会夸大平台价值，因此需要“撮合口径”与“资源配置物理口径”双轨计量。其三，本题无附件数据，需求、供给的时段分布须以官方统计与公开报道标定，标定结果必须通过基准情形与官方锚点（日成交约 130 万单、平均候车约 11.9 分钟）的一致性检验。方法路线上，将全日按北京市发改委高峰定义划分为早高峰、日间平峰、晚高峰、夜间四个时段，在均衡引擎基准输入下输出各时段成交、候车、未满足需求与空驶量，进而合成匹配度、打车难指数与需求加权综合匹配指数；空间维度以万人拥有量的功效系数做城际对比，以城市内三区供需比做内部结构刻画。

\subsection{补贴方案效果的反事实评价问题分析}

问题二的难点在于“补贴是否缓解打车难”表面上是数据问题，实质是识别问题：实际观测到的打车难变化，混杂了平台撮合技术带来的信息效应与补贴资金带来的激励效应，两者无法从前后对比中直接分离。本文的处理是构造两个反事实场景：无平台无补贴场景（S0）与平台无补贴场景（S1），两者之差即平台信息效应；各实际补贴方案（S2 至 S6）与 S1 之差即补贴资金效应。由此，评价问题转化为同一引擎在八种输入下的批量求解与成对比较。另一难点是参数不确定性的传播：需求弹性、供给弹性等关键参数取自文献区间，点估计下的结论可能对参数取值敏感，因此采用 2000 次蒙特卡洛仿真在参数区间内抽样，以方向概率和置信区间而非单点数值报告结论。识别框架还带来一个可检验的预测：若需求弹性大于供给弹性，乘客侧补贴刺激的需求增量会超过司机侧补贴增产的运力，双侧高额补贴在高峰时段可能反而加剧打车难，这一反直觉预测将由实际方案数据检验。

\subsection{预算约束下补贴结构的优化设计问题分析}

问题三是在问题二的评价框架上把补贴从外生参数变为内生决策变量：八维补贴向量（四个时段乘以乘客侧与司机侧）决定市场均衡，目标函数为全日失配成本（乘客未成交的福利损失加司机空驶的成本），约束为日补贴预算上限。难点有二。其一，目标函数含 $\min$、$\max$ 分段结构（成交取供求较小者、未满足需求与空驶取非负截断），且经不动点反馈后为非凸函数，不存在全局最优的解析充要条件，必须以“多起点一致性加网格下界验证”的证据链替代。其二，预算约束与物理上限（出租车总量管制下的出车率封顶、司机参与约束）使可行域边界复杂，需要逐约束标注活跃状态以解释最优解的结构。方法路线上，采用拉丁超立方多起点 SLSQP 求解，以全部起点收敛一致性作为数值最优性证据，以 39 万点双层网格枚举作为全局下界对照，并对预算规模做全区间扫描检验最优结构的稳健性。三问的关联是递进的：问题一提供基准确认打车难的时段位置，问题二识别出补贴的作用通道与反噬风险，问题三据此把资金集中配置到弹性最高的通道上。

% ================= 三、模型假设 =================
\section{模型假设}

\begin{enumerate}[label=\arabic*.~, leftmargin=2.5em, itemsep=2pt]
\item \textbf{封闭供给池假设}：出租车总量 $N_0=6.6$ 万辆在分析期内固定，平台与补贴只改变出车率，不吸引池外车辆进入。该假设使问题二的补贴效果估计为保守下界。
\item \textbf{时段独立性假设}：各时段内的需求与供给独立出清，忽略跨时段需求转移。灵敏度检验显示该简化对个别方案的高峰候车估计有最高 21.5\% 的偏差，结论方向不受影响，详见模型评价节。
\item \textbf{出车决策收入驱动假设}：司机在 $t$ 时段的出车数量只取决于该时段相对收入 $(p_0+s_d)/p_0$ 的弹性响应，出车率不超过 100\%（总量管制硬约束）。
\item \textbf{需求广义成本驱动假设}：乘客需求只随广义成本 $G=(p_0-s_p)+\rho_v W$ 变化，弹性 $\varepsilon_g$ 在时段间不变；$\varepsilon_g$ 为负，即广义成本上升需求下降。
\item \textbf{候车时间供需比律假设}：候车时间服从 $W=\hat{W}(S/D)^{-\sigma}$，供需比越低候车越长；$\sigma>0$ 取排队系统典型区间的中值偏保守水平。
\item \textbf{参数区间假设}：行为参数（$\varepsilon_g$、$\varepsilon_s$、$\sigma$、$\rho_v$）的真值落在文献报告区间内，区间内的参数不确定性由蒙特卡洛仿真传播，结论以方向概率报告。
\end{enumerate}

% ================= 四、符号说明 =================
\section{符号说明}

全文符号与求解代码变量名、术语表三方一致。决策变量见表 \ref{tab:dv}，主要参数见表 \ref{tab:param}，内生量与指标见表 \ref{tab:var}。

\begin{table}[H]
\centering
\caption{决策变量}
\label{tab:dv}
\zihao{-5}
\begin{tabular}{c c c}
\toprule
符号 & 说明 & 单位 \\
\midrule
$s_{d,t}$ & $t$ 时段司机端每单补贴（问题三决策变量，问题二为方案参数） & 元/单 \\
$s_{p,t}$ & $t$ 时段乘客端每单补贴（问题三决策变量，问题二为方案参数） & 元/单 \\
\bottomrule
\end{tabular}
\end{table}

\begin{table}[H]
\centering
\caption{主要参数}
\label{tab:param}
\zihao{-5}
\begin{tabular}{c c c}
\toprule
符号 & 说明 & 单位 \\
\midrule
$D_{0,t}$ & $t$ 时段基准需求（全日 155 万单按时段份额分摊的标定锚） & 万单 \\
$p_0$ & 基准运费（平均运距 8 公里约价） & 元/单 \\
$\varepsilon_g$ & 需求对广义成本的弹性，基准 $-0.62$，扫描区间 $[-1.2,-0.5]$ & -- \\
$\varepsilon_s$ & 司机供给对相对收入的响应弹性，基准 0.3，扫描 $[0.1,0.5]$ & -- \\
$\sigma$ & 候车时间对供需比的弹性，基准 0.8，扫描 $[0.5,1.5]$ & -- \\
$\rho_v$ & 乘客时间价值，基准 0.50，扫描 $[0.3,0.8]$ & 元/分钟 \\
$\hat{W}$ & 候车时间标定参数（使基准加权候车 11.9 分钟） & 分钟 \\
$N_0$ & 出租车总量（北京 6.6 万辆，行政管制上限） & 万辆 \\
$a_t$ & $t$ 时段基准出车率（标定，夜间 0.40 最低） & -- \\
$q_t$ & $t$ 时段单车服务率（时段时长除以平均服务时长） & 单/辆 \\
$\kappa$ & 平台信息效率乘数（无平台 1.0，有平台 1.15） & -- \\
$\varphi$ & 乘客失配单位社会成本（一单未成交的福利损失） & 元/单 \\
$\psi$ & 司机空驶单位成本（空驶间隔 3.38 公里乘变动成本） & 元/单 \\
$B$ & 平台日补贴预算，基准 100，扫描 $[0,300]$ & 万元/日 \\
\bottomrule
\end{tabular}
\end{table}

\begin{table}[H]
\centering
\caption{内生量与指标}
\label{tab:var}
\zihao{-5}
\begin{tabular}{c c c}
\toprule
符号 & 说明 & 单位 \\
\midrule
$D_t,\ S_t$ & $t$ 时段均衡需求与有效供给 & 万单 \\
$W_t,\ G_t$ & $t$ 时段均衡候车时间与乘客广义成本 & 分钟，元/单 \\
$Q_t,\ m_t$ & $t$ 时段成交单量 $\min(D_t,S_t)$ 与应答率 $\min(1,S_t/D_t)$ & 万单，-- \\
$U_t,\ E_t$ & $t$ 时段未满足需求与空驶过剩（空驶按物理口径计量） & 万单 \\
$M_t,\ H_t$ & $t$ 时段匹配度与打车难指数（物理口径） & -- \\
$SDMI$ & 综合供求匹配指数 $\sum_t\omega_t M_t$（$\omega_t$ 为需求份额权重） & -- \\
$Z$ & 全日失配成本 $\sum_t(\varphi U_t+\psi E_t)$，问题三目标函数 & 万元/日 \\
$B_{\mathrm{cost}}$ & 方案实际补贴支出 $\sum_t(s_{d,t}+s_{p,t})Q_t$ & 万元/日 \\
$\gamma,\ \eta$ & 万人拥有量与里程利用率（空间维度指标） & 辆/万人，-- \\
$\pi_t$ & $t$ 时段司机单位时间收入（参与约束校验用） & 元/小时 \\
$g$ & 均衡反馈增益 $\lvert\varepsilon_g\rvert\sigma\rho_v W/G$，$g<1$ 保证不动点唯一 & -- \\
\bottomrule
\end{tabular}
\end{table}

% ================= 五、模型的建立与求解 =================
\section{模型的建立与求解}

\subsection{分时段双边市场广义成本均衡引擎}

\subsubsection{问题结构分析}

出租车市场是典型的双边市场（Rochet \& Tirole, 2003; Armstrong, 2006），出租车行业的数量与价格管制及供求匹配问题亦有长期研究脉络（Douglas, 1972; Cairns \& Liston-Heyes, 1996; Yang \& Wong, 1998）：乘客侧的需求与司机侧的供给在同一时段内出清，候车时间把两侧行为耦合成反馈回路——需求越大，候车越长；候车越长，广义成本越高，需求又被抑制。平台的撮合技术与补贴资金都作用于这条回路。因此三问的共同内核是：给定时段、补贴与平台状态，求解该回路的均衡点。将全日按北京市发改委《关于调整本市出租汽车价格的通知》（京发改〔2013〕1052号）附件中的高峰定义划分为四个时段：早高峰（7:00-9:00，2 小时）、日间平峰（9:00-17:00，8 小时）、晚高峰（17:00-19:00，2 小时）、夜间（19:00-次日 7:00，12 小时）。时段之间独立出清（假设 2），每个时段的均衡由下列五个方程决定。

\subsubsection{模型建立}

\textbf{1. 需求侧：广义成本弹性需求。}乘客是否叫单，取决于一次打车的全部成本而非仅仅是运费。设 $t$ 时段候车时间为 $W_t$ 分钟，乘客时间价值为 $\rho_v$（元/分钟），乘客端补贴为 $s_{p,t}$（元/单），则乘客广义成本为
\begin{equation}
G_t = (p_0 - s_{p,t}) + \rho_v W_t . \label{eq:G}
\end{equation}
式 \eqref{eq:G} 的构成依据是运输经济学中广义出行成本的标准定义：货币支出加时间成本。以基准情形（无补贴均衡）的广义成本 $G_{0,t}$ 为参照，需求量按常弹性形式响应：
\begin{equation}
D_t = D_{0,t}\left(\frac{G_t}{G_{0,t}}\right)^{\varepsilon_g}, \qquad \varepsilon_g < 0 . \label{eq:D}
\end{equation}
常弹性形式的优点是无量纲、弹性可跨时段比较，且基准点处 $D_t=D_{0,t}$ 自动成立，便于标定。$\varepsilon_g$ 取 $-0.62$，由货币弹性（梁希喆基于北京 2001 至 2011 年数据的对数模型估计为 $-0.49$）与时间弹性（按文献区间上限取约 $-0.9$）按广义成本中的货币与时间份额合成，区间 $[-1.2,-0.5]$ 进入灵敏度与蒙特卡洛分析。

\textbf{2. 供给侧：收入弹性出车决策。}司机的出车决策看相对收入。设司机端补贴 $s_{d,t}$，则 $t$ 时段在线车辆数为
\begin{equation}
N_t = \min\left\{N_0\, a_t \left(\frac{p_0+s_{d,t}}{p_0}\right)^{\varepsilon_s},\ N_0\right\}, \label{eq:N}
\end{equation}
其中 $a_t$ 为基准出车率，$\varepsilon_s>0$ 为供给弹性（基准 0.3，取劳动供给弹性文献区间 0.1 至 0.5 的中值），$\min$ 截断对应出租车总量管制下出车率不超过 100\% 的物理约束（假设 3）。单车在 $t$ 时段可完成 $q_t$ 单（时段时长除以平均服务时长）。平台的作用以信息效率乘数 $\kappa$ 刻画：打车软件缩短司机找客的空驶间隔，同等车辆在同一时段可服务更多订单。由此区分两种口径的有效供给与物理供给：
\begin{equation}
S_t = \kappa N_t q_t \quad(\text{撮合口径}), \qquad S_t^{\mathrm{phys}} = N_t q_t \quad(\text{物理口径}) . \label{eq:S}
\end{equation}
双口径的纪律是：$\kappa$ 加快撮合周转，所以成交量、应答率、未满足需求、候车时间按撮合口径 $S_t$ 计算；$\kappa$ 不创造物理运力，所以空驶量与匹配度按物理口径 $S_t^{\mathrm{phys}}$ 计算。无平台时 $\kappa=1$，有平台时 $\kappa=1.15$（罗兰贝格报告空驶率由约 31\% 降至约 11\% 的保守折算）。

\textbf{3. 候车时间：供需比反馈律。}候车时间是市场拥挤程度的直接度量。供需比越低，乘客等车越久，设
\begin{equation}
W_t = \hat{W}\left(\frac{S_t}{D_t}\right)^{-\sigma}, \qquad \sigma > 0 , \label{eq:W}
\end{equation}
$\hat{W}$ 为标定参数（使基准加权候车与罗兰贝格报告一致），$\sigma$ 基准取 0.8（排队系统典型区间 0.5 至 1.5 的中值偏保守水平）。该形式满足两个必要条件：$S_t/D_t=1$ 时 $W_t=\hat{W}$；$S_t/D_t\to 0$ 时 $W_t\to\infty$，与排队系统趋于饱和时的行为一致。其与排队论模型（Gross \& Harris, 2008）的互证见问题一的模型验证。

\textbf{4. 均衡：不动点与唯一性。}把式 \eqref{eq:W} 代入式 \eqref{eq:G} 再代入式 \eqref{eq:D}，得到需求侧的合成映射
\begin{equation}
D_t = F(D_t) \equiv D_{0,t}\left[\frac{(p_0-s_{p,t}) + \rho_v \hat{W}\,(S_t/D_t)^{-\sigma}}{G_{0,t}}\right]^{\varepsilon_g} . \label{eq:fp}
\end{equation}
均衡需求是 $F$ 的不动点。对 $D_t$ 求导可得反馈增益
\begin{equation}
g \equiv \left|F'(D_t)\right| = \lvert\varepsilon_g\rvert\,\sigma\,\frac{\rho_v W_t}{G_t} . \label{eq:g}
\end{equation}
推导如下：$F$ 中随 $D_t$ 变化的部分是 $W_t=\hat{W}S_t^{-\sigma}D_t^{\sigma}$，故 $\mathrm{d}W_t/\mathrm{d}D_t=\sigma W_t/D_t$；由链式法则，$\mathrm{d}G_t/\mathrm{d}D_t=\rho_v\sigma W_t/D_t$；再对幂函数求导，$F'(D_t)=\varepsilon_g F(D_t)\cdot\rho_v\sigma W_t/(G_t D_t)$，在不动点 $F(D_t)=D_t$ 处约去 $D_t$ 即得式 \eqref{eq:g}。当 $g<1$ 时 $F$ 为压缩映像，由 Banach 不动点定理，均衡存在且唯一，简单迭代必收敛。基准参数下早高峰增益最大，$g_{\max}=0.62\times0.8\times0.5\times19.33/35.17=0.136<1$，四时段均满足压缩条件；蒙特卡洛抽样中出现的 $1\le g<1.8$ 的边角样本，由双初值（$0.5D_0$ 与 $2D_0$）收敛一致性确认唯一性，$g\ge1.8$ 的样本判为超出收敛域予以剔除（2000 次抽样中实际出现 0 次）。

\textbf{5. 市场出清、指标与成本分层。}均衡 $(D_t,S_t)$ 确定后，成交量、应答率、未满足需求与空驶量分别为
\begin{equation}
\begin{aligned}
Q_t&=\min(D_t,S_t),\qquad m_t=\min(1,S_t/D_t),\\
U_t&=\max(0,D_t-S_t),\qquad E_t=\max(0,S_t^{\mathrm{phys}}-Q_t) .
\end{aligned}\label{eq:clear}
\end{equation}
社会层面的失配成本按三个层次计量：乘客侧未成交的福利损失 $C_1=\sum_t\varphi U_t$（$\varphi$ 为单位失配成本，锚定平均运费 25.5 元与历史加价 5 至 20 元区间，取 20 元/单）；司机侧空驶的资源浪费 $C_2=\sum_t\psi E_t$（$\psi$ 为空驶单位成本，空驶间隔 3.38 公里乘变动成本 1.5 元/公里约 5 元/单）；平台侧资金消耗 $C_3=\sum_t(s_{d,t}+s_{p,t})Q_t$。问题三的目标函数取前两层的加总，即全日失配成本
\begin{equation}
Z = \sum_{t}\left(\varphi\,U_t + \psi\,E_t\right) . \label{eq:Z}
\end{equation}
$C_3$ 不进目标函数而作为预算约束，理由是：平台补贴是转移支付而非社会福利损失，但消耗平台资金，应在约束中体现。各时段失配成本以绝对量直接求和，不再乘需求权重，避免需求份额在 $D_0$ 标定与权重中的双重折算。

\subsubsection{模型求解：算法选择}

引擎的计算核心是式 \eqref{eq:fp} 的不动点求解。候选算法对比如表 \ref{tab:algo1} 所示。

\begin{table}[H]
\centering
\caption{不动点求解候选算法对比}
\label{tab:algo1}
\zihao{-5}
\begin{tabular}{p{3.4cm}p{4.6cm}p{5.2cm}}
\toprule
候选算法 & 优势 & 局限 \\
\midrule
阻尼不动点迭代 $D\leftarrow0.3D+0.7F(D)$ & 无需导数，负反馈结构下收敛域宽（$g<1.86$），实现简洁，逐时段独立 & 线性收敛，步数多于牛顿法 \\
牛顿法 & 二次收敛，步数少 & 需数值导数；分段结构（$\min/\max$ 截断）处导数不连续，可能越出收敛域 \\
解析近似（忽略反馈迭代一次） & 封闭形式 & 反馈恰是本题核心，近似误差不可控 \\
\bottomrule
\end{tabular}
\end{table}

选择阻尼不动点迭代，理由如下。第一，目标函数经 $\min/\max$ 截断后不可微，牛顿法的导数基础不稳固，而阻尼迭代只要求 $F$ 连续。第二，本题 $F$ 为减函数（需求越大候车越长、需求被抑制），负反馈下阻尼系数 0.7 的有效增益为 $|0.3-0.7g|$，在 $g<1.86$ 内均收敛，覆盖全部参数扫描区间；而简单迭代（阻尼系数 1）在 $g>1$ 时发散。第三，引擎需要在蒙特卡洛 2000 次抽样、八方案比较、22 个优化起点与 78 万点网格中被调用约百万次，单时段迭代平均约 10 步、最多 21 步（容差 $10^{-12}$），总耗时在秒级，满足批量调用需求。算法流程与三问在引擎上的运行方式见图 \ref{fig:flow}。

\begin{figure}[H]
\centering
\includegraphics[width=0.95\textwidth]{result_bundle/figures/fig_flowchart.pdf}
\caption{分时段双边市场均衡引擎与三问运行模式流程图}
\label{fig:flow}
\end{figure}

图 \ref{fig:flow} 展示了全文的计算主线。引擎左侧输入三类外生量：参数块（弹性 $\varepsilon_g$、$\varepsilon_s$、$\sigma$、时间价值 $\rho_v$ 与成本参数 $\varphi$、$\psi$）、标定块（基准需求 $D_0$、出车率 $a_t$、候车标定 $\hat{W}$、总量 $N_0$）与政策块（乘客补贴 $s_p$、司机补贴 $s_d$、平台乘数 $\kappa$）。引擎内核按时段执行四步：先由式 \eqref{eq:N}、\eqref{eq:S} 计算供给，再进入阻尼不动点回路求解需求，回路上方标注唯一性校验（增益 $g<1$ 用 Banach 判据，$1\le g<1.8$ 用双初值确认），最后按式 \eqref{eq:clear} 出清并输出指标。引擎右侧分岔出三问的运行模式：问题一固定政策块为基准与平台两档，读取指标；问题二让政策块取 S0 至 S6 八组输入，成对比较差值；问题三把政策块改为决策变量，外层套预算约束优化器。该图统摄全文：后续每一节都是图中某一条支路的展开，三问共享同一套方程与参数，保证了测量、评价、优化之间的口径一致性。

\subsubsection{基准标定与自检}

无附件数据下，标定的原则是“锚点外检”：先用公开数据确定参数，再让模型基准输出与未参与标定的官方锚点对比。标定输入包括：北京出租车总量 6.6 万辆（京政发〔2013〕13 号文件口径）、常住人口 2114.8 万（《北京市 2013 年国民经济和社会发展统计公报》）；基准运费 25.5 元/单（《法制晚报》2013 年报道的平均运距 8 公里约价）；平均候车 11.9 分钟、平台使空驶率由约 31\% 降至约 11\%（罗兰贝格研究报告）；全日成交 130 万车次以上（北京市交通委口径）。基准需求 $D_0$ 全日取 155 万单，即官方成交 130 万加上高峰未满足需求的估计；时段份额 $\{0.18, 0.42, 0.16, 0.24\}$ 与基准出车率 $\{0.75, 0.85, 0.72, 0.40\}$ 为构造标定（数据节如实声明），$\hat{W}=12.7$ 分钟使基准加权候车等于 11.9 分钟。基准自检结果见表 \ref{tab:calib}。

\begin{table}[H]
\centering
\caption{基准情形标定自检（模型输出与官方锚点对比）}
\label{tab:calib}
\zihao{-5}
\begin{tabular}{l c c l}
\toprule
检验量 & 模型基准值 & 官方锚点 & 判定 \\
\midrule
全日成交 $Q$（万单） & 134.6 & $\ge 130$ & 通过 \\
需求加权候车（分钟） & 11.85 & 11.9 & 通过（偏差 0.4\%） \\
早高峰应答率 $m$ & 0.591 & 高峰约六成量级 & 通过 \\
反馈增益 $g_{\max}$ & 0.136 & $<1$（唯一性条件） & 通过 \\
万人拥有量 $\gamma$（辆/万人） & 31.2 & 31.2（检索值） & 一致 \\
\bottomrule
\end{tabular}
\end{table}

五项检验全部通过，其中加权候车与全日成交两个外检锚点的偏差均在 1\% 以内，说明基准情形复现了 2014 年前后北京出租车市场的实际状态，后续三问的测量、评价与优化均建立在该通过外检的基准之上。

\subsection{问题一：时空双维度供求匹配测量}

\subsubsection{数据探索与初步判断}

对检索数据先做交叉观察，得到两个初步判断。第一，总量层面北京供给充裕：万人拥有量 $\gamma=31.2$ 辆/万人，超出住建部《城市道路交通规划设计规范》20 辆/万人的标准 56\%，在十个主要城市中最高；里程利用率 $\eta=68.8\%$（北京交通发展研究中心，2014 年 10 月），高于 65\% 的供求平衡经验阈值。若总量不足是打车难的主因，这两个数字不可能同时成立。第二，官方口径的全日成交（130 万单以上）与媒体报道的高峰长时间候车并存，提示失衡可能集中在特定时段与特定区域。据此形成待检验的工作假设：北京打车难的本质是时空结构错配而非总量短缺。该假设将由均衡引擎的分时段、分区域输出定量检验。需求时段份额、三区供求份额等无官方分项统计的量采用构造标定，其不确定性由蒙特卡洛仿真传播，并在结果节如实标注。

\subsubsection{指标体系建立}

时间维度构造四个原子指标与一个综合指数。匹配度
\begin{equation}
M_t=\frac{\min(S_t^{\mathrm{phys}},D_t)}{\max(S_t^{\mathrm{phys}},D_t)}\in[0,1],\qquad
H_t=\frac{U_t}{D_t}, \label{eq:MH}
\end{equation}
$M_t$ 对短缺与过剩对称惩罚，1 为完美匹配；$H_t$ 为打车难指数，即未满足需求占需求之比，直接刻画乘客侧困难程度；候车时间 $W_t$ 与应答率 $m_t$ 为可观测的过程指标。综合供求匹配指数取需求份额加权的时段平均：
\begin{equation}
SDMI = \sum_{t}\omega_t M_t,\qquad \omega_t = D_{0,t}\Big/\sum_{\tau}D_{0,\tau} . \label{eq:sdmi}
\end{equation}
权重只用于指标合成，不进入问题三的目标函数。空间维度分两层：城际层用万人拥有量的功效系数 $\gamma/\gamma_{\mathrm{std}}$（$\gamma_{\mathrm{std}}=20$ 为标准值，系数大于 1 即达标）做十城对比；城内层把城市划分为核心区、居住区、郊区三类功能分区，以各分区供需比刻画内部错配。平台前后的对比按双口径纪律执行：撮合口径指标（$Q,m,U,W$）含 $\kappa$，物理口径指标（$E,M,SDMI$）不含 $\kappa$。

\subsubsection{模型求解}

让引擎以测量模式运行两次：第一次 $\kappa=1$（传统扬招），第二次 $\kappa=1.15$（平台撮合），政策块补贴均为零。两次运行共用同一组标定参数，差别只在供给方程中的信息效率乘数，因此输出之差可全部归因于平台撮合技术。两口径八次时段求解的不动点迭代平均 10.5 步、最多 21 步收敛（扬招口径迭代初值即基准均衡，一步收敛），总耗时不足 0.1 秒。

\subsubsection{结果与结果分析}

时间维度的双口径测量结果见表 \ref{tab:q1t} 与图 \ref{fig:q1t}。

\begin{table}[H]
\centering
\caption{问题一时间维度测量结果（北京，扬招 $\kappa=1$ 对平台 $\kappa=1.15$）}
\label{tab:q1t}
\zihao{-5}
\setlength{\tabcolsep}{4pt}
\begin{tabular}{l c c c c c c c}
\toprule
时段 & $D$（万单） & $S$（万单） & $W$（分钟） & $m$ & $U$（万单） & $E$（万单） & $M$ \\
\midrule
早高峰（扬招） & 27.9 & 16.50 & 19.33 & 0.591 & 11.40 & 0 & 0.591 \\
早高峰（平台） & 28.4 & 18.97 & 17.51 & 0.669 & 9.38 & 0 & 0.582 \\
日间平峰（扬招） & 65.1 & 99.73 & 9.03 & 1.000 & 0 & 34.63 & 0.653 \\
日间平峰（平台） & 65.7 & 114.69 & 8.13 & 1.000 & 0 & 34.02 & 0.659 \\
晚高峰（扬招） & 24.8 & 15.84 & 18.18 & 0.639 & 8.96 & 0 & 0.639 \\
晚高峰（平台） & 25.2 & 18.22 & 16.46 & 0.723 & 6.97 & 0 & 0.629 \\
夜间（扬招） & 37.2 & 79.20 & 6.94 & 1.000 & 0 & 42.00 & 0.470 \\
夜间（平台） & 37.5 & 91.08 & 6.24 & 1.000 & 0 & 41.72 & 0.473 \\
\midrule
全日加权/加总（扬招） & \multicolumn{2}{c}{$Q=134.64$ 万单} & 11.85 & \multicolumn{2}{c}{$SDMI=0.5955$} & \multicolumn{2}{c}{$Z=790.4$ 万元/日} \\
全日加权/加总（平台） & \multicolumn{2}{c}{$Q=140.38$ 万单} & 10.70 & \multicolumn{2}{c}{$SDMI=0.5956$} & \multicolumn{2}{c}{$Z=705.8$ 万元/日} \\
\bottomrule
\end{tabular}
\end{table}

\begin{figure}[H]
\centering
\includegraphics[width=0.95\textwidth]{result_bundle/figures/fig_q1_temporal.pdf}
\caption{不同时段出租车供求匹配程度（北京，扬招与平台双口径对比）}
\label{fig:q1t}
\end{figure}

表 \ref{tab:q1t} 与图 \ref{fig:q1t} 给出的图景高度一致：北京的供求失衡沿时段轴呈“两端短缺、中间过剩”的锯齿结构。早高峰需求 27.9 万单对供给 16.5 万单，缺口 11.4 万单，应答率 0.591，候车 19.33 分钟，打车难指数 0.409，是全日最紧张的时段；晚高峰缺口 8.96 万单，候车 18.18 分钟，紧张程度次之；日间平峰与夜间则反转，空驶当量分别达到 34.63 万单与 42.0 万单，夜间匹配度仅 0.470，大量运力闲置。

图中平台档（深色柱）相对扬招档（浅色柱）的变化揭示了撮合技术的作用边界：早高峰候车下降 1.82 分钟、未满足需求减少 2.02 万单、应答率上升 7.8 个百分点，但匹配度 $M$ 不升反微降（0.591 到 0.582），因为平台没有增加一辆车，只是让既有车辆周转更快，成交的 5.74 万单增量（140.38 对 134.64）全部来自周转加速。全日汇总行中，加权候车从 11.85 分钟降至 10.70 分钟，失配成本从 790.4 万元降至 705.8 万元，而 SDMI 几乎不动（0.5955 对 0.5956）：物理资源配置未变，撮合效率改善了乘客体验。这一“SDMI 不变、打车难缓解”的分离正是双口径计量的信息量所在，也为问题二分解信息效应与资金效应埋下伏笔。

空间维度的测量结果见图 \ref{fig:q1s}。

\begin{figure}[H]
\centering
\includegraphics[width=0.95\textwidth]{result_bundle/figures/fig_q1_spatial.pdf}
\caption{空间维度供求匹配：十城万人拥有量功效系数与北京城内三区供需比}
\label{fig:q1s}
\end{figure}

图 \ref{fig:q1s} 左面板按万人拥有量功效系数排序给出十个城市：北京以 31.2 辆/万人居首，功效系数 1.56；上海、天津在 21 辆/万人附近勉强达标；广州、深圳、成都、重庆等七城在 6.7 至 17.1 辆/万人之间，均未达到 20 辆/万人的标准。北京的总量充裕度在全国领先，这与“北京打车全国闻名难打”的现实形成鲜明反差，说明总量指标不能解释打车难的城际差异。右面板转向北京内部：核心区供需比 1.222，车辆相对需求过剩；居住区供需比 0.857，早高峰大量通勤需求从居住区涌向核心区而车辆滞留在核心区；郊区供需比 0.750，缺口最大。三区供需比从 1.22 到 0.75 的梯度证明，即便在总量超标的北京，供给在空间上的分布仍然与需求错位：乘客在居住区和郊区等车，车辆在核心区空驶。时间维度“高峰短缺、平峰过剩”与空间维度“外围短缺、核心过剩”叠加，共同回答了问题一：北京出租车资源在总量上是充裕的（万人拥有量超标准 56\%），打车难来自时空双维度的结构性错配，高峰时段与外围区域是短缺的位置。这一判断的政策含义直接：增加总量的边际收益很低，能改变时空分布或撮合效率的手段才是有效杠杆。

\subsubsection{结果验证与模型验证}

结果验证（最优性维度）：问题一为测量任务，无优化目标，其“最优性”等价于均衡解的数值可靠性。四时段不动点迭代均在 21 步内收敛至 $10^{-12}$ 容差，反馈增益 $g_{\max}=0.136<1$ 满足 Banach 压缩条件，均衡唯一性有理论保证；基准输出通过表 \ref{tab:calib} 的五项外检锚点检验。

模型验证（独立结构互证维度）：引入与均衡模型结构独立的 M/M/c 排队模型交叉验证。以时段到达率 $\lambda=D_t/$时段时长、单车服务率 $\mu=1/$平均服务时长、服务台数 $c=N_t$ 计算 Erlang-C 平均排队等待 $W_q$ 与系统利用率 $\rho=\lambda/(c\mu)$，结果见图 \ref{fig:valid} 左面板。

\begin{figure}[H]
\centering
\includegraphics[width=0.95\textwidth]{result_bundle/figures/fig_validation.pdf}
\caption{模型验证：M/M/c 排队论独立互验（左）与蒙特卡洛批次收敛（右）}
\label{fig:valid}
\end{figure}

图 \ref{fig:valid} 左面板显示，扬招基准下早高峰 $\rho=1.69$、晚高峰 $\rho=1.57$，均超过 1，排队系统无稳态、等待时间发散。这不是模型误差的证据，恰恰相反：两个结构完全不同的模型对同一时段给出同方向的判断——均衡模型算出高峰候车 18 至 19 分钟量级，排队模型直接判定系统过载，高峰打车难获得了独立的结构性证据。日间平峰与夜间 $\rho$ 分别为 0.65 与 0.47，排队模型给出 $W_q\approx0$：数万辆服务台组成的大服务台池在利用率小于 1 时，等待概率按指数趋零（大偏差理论），拥挤排队分量可忽略；此时均衡模型给出的 6 至 9 分钟候车主要由物理接驾分量构成（空驶间隔 3.38 公里对应的行驶时间）。两个模型对照完成了候车时间的两分量分解：候车总时间等于拥挤排队分量加物理接驾分量，高峰以前者为主，平峰以后者为主。图 \ref{fig:valid} 右面板为问题二蒙特卡洛仿真的批次均值收敛曲线，三条效应估计序列在 500 次后均进入平稳带，2000 次抽样的均值标准差与均值之比不超过 1.2\%，表明仿真次数充分、均值估计稳定，其数值细节在问题二节展开。

\subsection{问题二：补贴方案的反事实评价}

\subsubsection{模型建立：反事实识别策略}

评价“补贴是否缓解打车难”的障碍是效应混杂：补贴方案上线前后的任何变化，都同时包含平台撮合技术（信息效应 $\kappa$）与补贴资金（激励效应 $s_p,s_d$）的贡献。本文构造八个场景把两者拆开。设无平台无补贴场景 S0（$\kappa=1$，零补贴）与平台无补贴场景 S1（$\kappa=1.15$，零补贴），则信息效应定义为
\begin{equation}
\Delta_{\mathrm{info}} = Y(\mathrm{S1}) - Y(\mathrm{S0}), \label{eq:info}
\end{equation}
$Y$ 为任一结果指标。五个实际补贴方案 S2 至 S6 与 S1 之差即各方案的纯资金效应：
\begin{equation}
\Delta_{\mathrm{sub}}^{(k)} = Y(\mathrm{S}k) - Y(\mathrm{S1}),\qquad k=2,\dots,6 . \label{eq:sub}
\end{equation}
实际方案相对“没有平台的时代”的净效果为两项之和：$Y(\mathrm{S}k)-Y(\mathrm{S0})=\Delta_{\mathrm{info}}+\Delta_{\mathrm{sub}}^{(k)}$。六个实际补贴场景的参数取自 2014 年补贴大战的公开报道，见表 \ref{tab:schemes}。

\begin{table}[H]
\centering
\caption{反事实场景与实际补贴方案参数（元/单）}
\label{tab:schemes}
\zihao{-5}
\begin{tabular}{l l l c c}
\toprule
场景 & 方案背景 & 时间 & 乘客端 $s_p$ & 司机端 $s_d$ \\
\midrule
S0 & 无平台无补贴（扬招基准） & -- & 0 & 0 \\
S1 & 平台无补贴（信息效应反事实） & -- & 0 & 0 \\
S2 & 滴滴“乘客 10 元加司机 10 元” & 2014-01 & 10 & 10 \\
S3 & 快的“乘客 10 元加司机 15 元” & 2014-01 & 10 & 15 \\
S4 & 滴滴补贴下调“乘客 5 元加司机 10 元” & 2014-02-10 & 5 & 10 \\
S5 & 快的北京“乘客 11 元加司机高峰 11/其他 5” & 2014-02-18 & 11 & 高峰 11，其他 5 \\
S6 & 补贴大战退潮后暗补期“乘客 0 加司机约 3” & 2014-05 后 & 0 & 约 3 \\
\bottomrule
\end{tabular}
\end{table}

\subsubsection{模型求解：蒙特卡洛仿真设计}

行为参数的真值只知道区间，点估计下的评价结论可能系于参数取值的运气。为此在参数区间内做蒙特卡洛仿真：$\varepsilon_g$ 在 $[-1.2,-0.5]$、$\varepsilon_s$ 在 $[0.1,0.5]$、$\sigma$ 在 $[0.5,1.5]$、$\rho_v$ 在 $[0.3,0.8]$ 上独立均匀抽样，每次抽样求解全部八个场景的均衡并记录效应差值，共 2000 次（固定种子 42）。报告口径为蒙特卡洛均值与 5\% 至 95\% 分位区间，方向结论以方向概率（效应为正/负的样本比例）报告。点值（参数取基准值时的单次结果）只作对照，因为效应函数对参数非线性，$E[f(X)]\ne f(E[X])$，均值与点值的差异是 Jensen 不等式的正常结果而非模型错误。

\subsubsection{结果与结果分析}

八场景的核心结果见表 \ref{tab:q2r} 与图 \ref{fig:q2s}。

\begin{table}[H]
\centering
\caption{问题二各场景评价结果（高峰为早晚高峰均值/合计，$Z$ 为全日失配成本）}
\label{tab:q2r}
\zihao{-5}
\setlength{\tabcolsep}{3pt}
\begin{tabular}{l c c c c c c}
\toprule
场景 & $W$（分钟） & $H$ & $U$（万单） & $Z$（万元/日） & 补贴支出（万元） & 效益比 \\
\midrule
S0 无平台 & 18.76 & 0.385 & 20.36 & 790.4 & 0 & -- \\
S1 平台无补贴 & 16.99 & 0.304 & 16.36 & 705.8 & 0 & -- \\
S2 滴滴 2014-01 & 18.39 & 0.369 & 24.21 & 812.8 & 3458.7 & $-0.006$ \\
S3 快的 2014-01 & 17.90 & 0.348 & 22.95 & 825.1 & 4377.1 & $-0.008$ \\
S4 滴滴 2014-02 & 16.98 & 0.304 & 18.03 & 769.7 & 2352.9 & $+0.009$ \\
S5 快的 2014-02 & 18.60 & 0.379 & 25.37 & 775.8 & 3073.2 & $+0.005$ \\
S6 暗补期 2014-05 后 & 16.59 & 0.283 & 15.29 & 713.9 & 425.5 & $+0.180$ \\
\bottomrule
\end{tabular}
\end{table}

\begin{figure}[H]
\centering
\includegraphics[width=0.95\textwidth]{result_bundle/figures/fig_q2_schemes.pdf}
\caption{各补贴方案反事实评价（S0 无平台基准至 S6 暗补期）}
\label{fig:q2s}
\end{figure}

表 \ref{tab:q2r} 与图 \ref{fig:q2s} 揭示了三个层次的事实。第一层，平台本身是有价值的：S1 相对 S0，高峰候车从 18.76 分钟降到 16.99 分钟，高峰未满足需求减少 4.00 万单，日失配成本下降 84.5 万元，这是没有任何补贴时撮合技术带来的纯信息效应。

第二层，补贴大战的主力方案是负效果：S2、S3 的高峰未满足需求不降反升，分别比 S1 多 7.85 万单和 6.59 万单，日失配成本分别恶化 107.0 万元和 119.3 万元，而每日补贴支出高达 3458.7 万元和 4377.1 万元；即便与没有平台的时代（S0）相比，S2 的净效果仍是高峰未满足需求增加 3.85 万单、日失配成本恶化 22.5 万元，信息效应的正贡献被资金效应的负贡献吞没。

第三层，补贴结构比补贴规模重要：S6 只给司机侧约 3 元的小额补贴，日支出仅 425.5 万元，却是六个方案中唯一高峰未满足需求低于 S1 的方案（15.29 万单），日失配成本比 S0 低 76.5 万元，成本效益比 0.18，比 S4 的 0.009 高一个数量级。图 \ref{fig:q2s} 中三个面板的对比直观显示：补贴支出柱（S2、S3 最高）与失配成本柱的排序几乎相反，烧钱最多的方案恰恰是最差的方案。

效应分解的定量结果见图 \ref{fig:q2d}。

\begin{figure}[H]
\centering
\includegraphics[width=0.9\textwidth]{result_bundle/figures/fig_q2_decompose.pdf}
\caption{信息效应与资金效应分解（高峰候车时间，基准参数点值；蒙特卡洛均值与 90\% 区间见正文及图 \ref{fig:valid}）}
\label{fig:q2d}
\end{figure}

图 \ref{fig:q2d} 以瀑布形式展示高峰候车时间从 S0 到 S2 的分解，图中柱高为基准参数点值：信息效应（S1 减 S0）$-1.77$ 分钟、资金效应（S2 减 S1）$+1.40$ 分钟，S0 的 18.76 分钟经两步变化到 S2 的 18.39 分钟。两效应的不确定性由 2000 次蒙特卡洛仿真刻画：信息效应的均值为 $-2.23$ 分钟，90\% 区间 $[-3.56,-1.08]$，2000 个样本全部为负，方向概率 100\%，平台撮合技术缩短高峰候车的结论高度稳健；点值 $-1.77$ 分钟落在区间之内，与均值的偏差 26.2\% 源于效应函数对需求弹性的凸性（弹性越负，需求回流越弱，候车改善越大）。资金效应的均值为 $+2.69$ 分钟，90\% 区间 $[0.72,5.34]$，方向概率 99.65\%：滴滴 2014 年 1 月方案使高峰候车显著回升，点值 $+1.40$ 分钟处于分布低端（偏差 92.3\%，反噬幅度随需求弹性绝对值非线性放大）。两个效应叠加后，S2 相对 S0 的净效果接近零甚至为负。

这一反直觉结果的机理清晰可查：乘客端补贴 10 元直接压低广义成本，按需求弹性 0.62 放大需求；司机端补贴 10 元按供给弹性 0.3 增产出车。需求侧杠杆的强度约为供给侧的两倍，双侧等额补贴造成的需求增量超过运力增量，高峰时段的供需缺口被拉大而非缩小。需要声明的是，该结论依赖弹性序 $\lvert\varepsilon_g\rvert>\varepsilon_s$：若真实弹性序反转（供给弹性大于需求弹性绝对值），资金效应方向将随之反转，此时的最优结构应转向乘客侧；本文弹性取值的文献依据与区间扫描（$\varepsilon_g$ 至 $-0.5$、$\varepsilon_s$ 至 $0.5$ 的边角组合）见灵敏度分析节，在扫描区间内反噬方向不改变。

\subsubsection{结果验证与模型验证}

结果验证（数值可靠性维度）：2000 次抽样全部落在收敛域内（$g\ge1.8$ 剔除率为 0），图 \ref{fig:valid} 右面板显示信息效应、资金效应与 S6 对 S2 成本差三条批次均值序列在 500 次后进入平稳带，均值标准差与均值之比分别为 0.78\%、1.20\%、0.96\%，仿真收敛充分。

模型验证（不确定性传播维度）：三条核心判断在 2000 次仿真中分别以 100\%、99.65\%、99.75\% 的样本成立——平台信息效应使高峰候车下降、S2 资金效应使高峰候车上升、S6 优于 S2。结论对参数区间内的不确定性稳健。另有两项口径说明：其一，空驶率由 29.7\% 降至 26.9\% 的幅度小于公开报道“31\% 降至约 11\%”的叙事，两者口径不同，前者为本文的时段加权空驶间隔口径，后者为运营统计口径，$\kappa$ 的取值已按保守方向折算；其二，模型为封闭供给池，未含补贴大战期间社会车辆涌入带来的供给扩大，故问题二的补贴效果估计为保守下界，供给池扩大只会增加运力、强化上述结论。

\subsection{问题三：预算约束下的补贴结构优化}

\subsubsection{模型建立}

新平台的决策是八维补贴向量 $\boldsymbol{x}=(s_{d,1},\dots,s_{d,4},s_{p,1},\dots,s_{p,4})$，目标为全日失配成本最小，约束为日补贴预算：
\begin{equation}
\min_{\boldsymbol{x}}\ Z(\boldsymbol{x})=\sum_t\left(\varphi U_t(\boldsymbol{x})+\psi E_t(\boldsymbol{x})\right) \label{eq:q3obj}
\end{equation}
\begin{equation}
\text{s.t.}\quad B_{\mathrm{cost}}(\boldsymbol{x})=\sum_t (s_{d,t}+s_{p,t})\,Q_t(\boldsymbol{x}) \le B,\qquad 0\le s_{d,t},\,s_{p,t}\le 20 . \label{eq:q3con}
\end{equation}
其中 $U_t,E_t,Q_t$ 由均衡引擎式 \eqref{eq:fp}、\eqref{eq:clear} 内生决定，$\kappa=1.15$（新平台具备撮合技术）。基准预算 $B=100$ 万元/日，按滴滴 2014 年三个月约 14 亿元全国补贴（亿邦动力网 2014-06-10 报道）乘以北京约 10\% 份额折算；单补贴上限 20 元/单对标补贴大战峰值区间。目标函数含 $\min/\max$ 分段结构并经不动点反馈，为非凸函数，全局最优无解析充要条件，最优性须由数值证据链支持。问题二的经验直接提示了结构猜测：资金应投向弹性更高的一侧（司机侧）与缺口更大的时段（高峰），优化将检验这一猜测并给出定量配置。

\subsubsection{模型求解：算法选择与最优性证据链}

候选算法对比见表 \ref{tab:algo3}。

\begin{table}[H]
\centering
\caption{问题三优化候选算法对比}
\label{tab:algo3}
\zihao{-5}
\begin{tabular}{p{3.6cm}p{4.6cm}p{5.0cm}}
\toprule
候选算法 & 优势 & 局限 \\
\midrule
多起点 SLSQP & 处理预算约束与非负边界成熟，每次迭代调用引擎代价低 & 局部方法，非凸问题需多起点与下界对照 \\
网格枚举 & 提供全局下界，结果直观 & 维数灾难：八维每维 5 档即 39 万点，只能做粗下界 \\
动态规划 & 序贯决策问题最优 & 本题为单期配置问题，无阶段递推结构，不适用 \\
进化算法 & 不依赖梯度，全局探索 & 预算约束需罚函数处理，精度与收敛速度低于 SLSQP \\
\bottomrule
\end{tabular}
\end{table}

采用“多起点 SLSQP 主求解加网格枚举下界验证”的组合方案：以 20 个拉丁超立方随机起点加 2 个结构起点（纯司机侧、纯乘客侧）运行 SLSQP（最大迭代 300，函数容差 $10^{-10}$），若全部起点收敛到同一解，则该解是强局部最优的数值证据；再以粗网格（每维 $\{0,5,10,15,20\}$，39.06 万点）与细网格（最优解邻域加密，39.06 万点）枚举给出全局下界，若 SLSQP 解不劣于网格最优且差距在容差内，则全局最优性获得双证据支持。预算维度另做 $B\in[0,300]$ 万元的 13 点全扫描，检验最优结构对预算规模的稳健性。

\subsubsection{结果与结果分析}

优化结果见表 \ref{tab:q3r} 与图 \ref{fig:q3o}。

\begin{table}[H]
\centering
\caption{问题三最优补贴方案（$B=100$ 万元/日）}
\label{tab:q3r}
\zihao{-5}
\setlength{\tabcolsep}{4pt}
\begin{tabular}{l c c c c}
\toprule
时段 & 司机端 $s_d$（元/单） & 乘客端 $s_p$（元/单） & 支出（万元/日） & 约束状态 \\
\midrule
早高峰 & 2.36 & 0 & 46.0 & $s_p$ 下界活跃 \\
日间平峰 & 0 & 0 & 0 & 双侧下界活跃，参与约束活跃 \\
晚高峰 & 2.87 & 0 & 54.0 & $s_p$ 下界活跃 \\
夜间 & 0 & 0 & 0 & 双侧下界活跃，参与约束活跃 \\
\midrule
合计 & \multicolumn{2}{c}{预算约束活跃（支出 100.0 万元/日）} & 100.0 & 总量上限四时段均松弛 \\
\bottomrule
\end{tabular}
\end{table}

\begin{figure}[H]
\centering
\includegraphics[width=0.95\textwidth]{result_bundle/figures/fig_q3_optimal.pdf}
\caption{最优补贴结构与高峰改善效果}
\label{fig:q3o}
\end{figure}

表 \ref{tab:q3r} 与图 \ref{fig:q3o} 显示，最优方案呈极端结构：100 万元预算全部投向司机侧（乘客侧补贴四时段均为零，下界约束活跃）、全部投向早晚高峰（平峰与夜间补贴为零），早高峰 2.36 元/单、晚高峰 2.87 元/单，支出按 46.0 比 54.0 分配。这一结构是问题二机理的镜像：司机侧补贴的杠杆（供给弹性通道）不会被需求侧的反噬抵消，高峰时段的边际失配成本最高，资金自然向“司机侧乘高峰”的通道集中。平峰与夜间的双侧下界活跃配合司机参与约束活跃，说明这些时段司机收入已处于出车意愿的临界水平，但补贴投在这里的边际收益仍低于高峰。方案效果上，日失配成本从平台无补贴情形的 705.8 万元降至 687.2 万元，改善 2.64\%；高峰加权候车下降 2.1\%，高峰打车难指数下降 6.1\%，早高峰打车难指数 0.315、晚高峰 0.256，各时段乘客广义成本均不高于无补贴基准，公平性约束（$G_t\le G_{0,t}$ 全时段成立）自动满足。需要强调的定位是：2.64\% 这个点值对基准需求水平敏感（$D_0$ 扰动 20\% 时 $Z^*$ 摆动约 33\%），真正稳健的结论是最优结构而非改善幅度。

预算扫描结果见图 \ref{fig:q3p}。

\begin{figure}[H]
\centering
\includegraphics[width=0.95\textwidth]{result_bundle/figures/fig_q3_pareto.pdf}
\caption{补贴预算扫描：权衡曲线、边际收益与最优结构稳定性}
\label{fig:q3p}
\end{figure}

图 \ref{fig:q3p} 左面板给出 $B$ 从 0 到 300 万元的权衡曲线：日失配成本从 705.8 万元单调降至 655.2 万元，曲线全程下凸，边际收益（每万元补贴换取的失配成本下降）从 0.196 平滑递减至 0.148，没有出现突变拐点。右面板显示乘客侧资金占比与高峰资金占比两条线：在 12 个正预算点上，前者恒为 0、后者恒为 100\%（零预算点无补贴，占比无定义），最优结构在 $[0,300]$ 万元全区间不发生转换。两条曲线合起来的结论是：无论平台资金实力如何，“纯司机侧加纯高峰”都是最优配置，预算只改变补贴强度（早高峰补贴从 0.30 元/单增至 7.19 元/单），不改变补贴结构。这一负结果（无结构转换点）本身就是重要的管理启示：平台不需要随预算增长重新设计补贴方案，结构设计的决策可以一次性完成。

\subsubsection{结果验证}

最优性证据链有三环。第一环，多起点一致性：22 个起点（20 个拉丁超立方加 2 个结构起点）全部收敛到同一解，目标值相对标准差 $4.9\times10^{-14}$，数值上不可区分。第二环，网格下界：粗网格 39.06 万点的最优 $Z=688.6$ 万元，SLSQP 解 687.2 万元优于粗网格下界方向（相对差距 0.21\%）；细网格在最优解邻域内 39.06 万点的最优 $Z=688.0$ 万元，相对差距 0.11\%，两层网格均确认 SLSQP 解不劣于枚举结果。第三环，结构稳健性：$\varepsilon_g$ 取 $-0.9$ 的强需求弹性复核，最优方向不变；$\varphi/\psi$ 比值在 25 组网格组合上的扫描显示，比值不超过 2.00 时最优结构翻转为乘客侧主导，不低于 2.50 时维持司机侧主导，临界区间 $(2.0,2.5)$，而基准比值 4.0 距离临界值的安全边际超过 60\%，常规参数扰动区间 $[3.2,4.8]$ 完全落在稳健域内。三环证据共同支持结论：最优解是数值可信的强局部最优，且其结构结论在参数合理扰动下不变。模型求解全程（22 起点优化加 78 万点网格加 13 点预算扫描）耗时 57.4 秒，引擎批量调用的效率满足验证需求。

% ================= 六、灵敏度与误差分析 =================
\section{灵敏度与误差分析}

\subsection{灵敏度分析}

对八个关键参数做 $\pm 20\%$ 单参数扰动（与求解代码的扰动范围一致），分别考察平台无补贴情形的失配成本 $Z$ 与重优化后的最优失配成本 $Z^*$，结果见图 \ref{fig:tornado} 与表 \ref{tab:tornado}。

\begin{figure}[H]
\centering
\includegraphics[width=0.95\textwidth]{result_bundle/figures/fig_sensitivity_tornado.pdf}
\caption{Tornado 灵敏度排序：平台无补贴失配成本 $Z$（左）与重优化最优值 $Z^*$（右），扰动范围 $\pm 20\%$}
\label{fig:tornado}
\end{figure}

\begin{table}[H]
\centering
\caption{参数灵敏度排序（$\pm 20\%$ 扰动下目标值的相对摆幅）}
\label{tab:tornado}
\zihao{-5}
\begin{tabular}{l c c c c}
\toprule
参数 & 基准值 & $Z$ 摆幅（平台无补贴） & $Z^*$ 摆幅（重优化） & 排序 \\
\midrule
$D_0$ 基准需求 & 155 万单/日 & 31.8\% & 33.0\% & 1 \\
$\psi$ 空驶单位成本 & 5 元/单 & 21.5\% & 22.0\% & 2 \\
$\varphi$ 失配单位成本 & 20 元/单 & 18.5\% & 18.0\% & 3 \\
$\sigma$ 候车弹性 & 0.8 & 0.94\% & 1.18\% & 4 \\
$\varepsilon_s$ 供给弹性 & 0.3 & 0.00\% & 1.08\% & 5 \\
$\varepsilon_g$ 需求弹性 & $-0.62$ & 0.62\% & 0.82\% & 6 \\
$\hat{W}$ 候车标定 & 12.7 分钟 & 0.43\% & 0.58\% & 7 \\
$\rho_v$ 时间价值 & 0.5 元/分钟 & 0.43\% & 0.58\% & 8 \\
\bottomrule
\end{tabular}
\end{table}

图 \ref{fig:tornado} 的两个面板呈现一致的梯队结构。第一梯队是 $D_0$、$\psi$、$\varphi$ 三个总量与成本参数：$D_0$ 扰动 20\% 引起 $Z$ 摆动约 32\%，因为失配成本近似与市场规模成正比；$\psi$ 与 $\varphi$ 的摆幅在 18\% 至 22\%，它们直接是目标函数的价格权重。第二梯队是弹性与行为参数，摆幅全部不超过 1.2\%：$\sigma$、$\varepsilon_g$、$\rho_v$、$\hat{W}$ 只影响均衡在反馈回路上的落点，对全日加总成本的影响被时段间的抵消效应摊薄。左面板中 $\varepsilon_s$ 的摆幅恰为零，这不是计算错误而是结构事实：平台无补贴情形 $s_d=0$，供给方程中的弹性因子 $((p_0+0)/p_0)^{\varepsilon_s}=1$ 恒成立，供给对弹性免疫；右面板重优化后 $\varepsilon_s$ 摆幅升至 1.08\%，补贴通过该弹性起作用，排序从末位升到第五位。两个面板对照给出一个明确结论：模型的绝对成本水平由 $D_0$、$\psi$、$\varphi$ 主导（这正是点值结论不宜过度解读的原因），而弹性族参数的低灵敏度保证了问题二、问题三的核心判断（信息效应为正、补贴反噬、司机侧高峰结构）对弹性取值不敏感。

\subsection{策略转换临界点}

最优补贴结构的稳健性以两类临界点刻画。第一类沿 $\varphi/\psi$ 维度：在 $\varphi\in\{10,15,20,25,30\}$ 与 $\psi\in\{2,3.5,5,7.5,10\}$ 的 25 组网格组合上重优化，6 组发生结构翻转，全部位于比值 2.0 及以下；比值不低于 2.50 的 19 组全部维持司机侧主导。临界比值落在区间 $(2.0,2.5)$ 内，基准比值 $20/5=4.0$ 距上界的安全边际为 $(4.0-2.5)/2.5=60\%$，$\pm 20\%$ 扰动给出的比值区间 $[3.2,4.8]$ 与翻转域无交集。经济解释是：$\varphi/\psi$ 度量乘客失配相对司机空驶的“价格”，该比值越高，缓解乘客侧短缺越值钱，资金越应投向能增产运力的司机侧；当空驶成本相对升高（比值跌破 2.0），乘客侧降价压需求、减少空驶占用的组合才变得划算。第二类沿预算维度：$B\in[0,300]$ 万元 13 点扫描未发现结构转换点（负结果），边际收益从 0.196 到 0.148 平滑递减，最优结构对预算规模稳健。

\subsection{活跃约束标注}

问题三最优解处各约束的状态在结果节表 \ref{tab:q3r} 已逐条标注，此处汇总其经济学含义。预算约束活跃：100 万元全部用尽，补贴的边际收益（0.178）仍为正，资金是稀缺资源。乘客侧补贴四个时段下界活跃：放松零补贴限制不会改善目标，乘客侧通道在当前弹性序下无投资价值。平峰与夜间司机补贴下界活跃且司机参与约束活跃：这些时段司机收入恰在出车意愿临界水平，补贴若投向此处只能覆盖参与约束的固定门槛，边际收益低于高峰。出租车总量上限四时段均松弛：补贴引致的出车率提升最高约 3\%，远未触及 6.6 万辆的管制天花板，总量管制不是优化解的束缚，这与问题一“总量充裕、结构错配”的判断自洽。

\subsection{误差来源、量化与可接受性}

本文误差按来源分为四类，量化与处置如下。

\textbf{1. 参数标定误差。}弹性与行为参数取自文献区间而非本题实测，是最大误差源。处置方式是区间化：$\varepsilon_g\in[-1.2,-0.5]$、$\varepsilon_s\in[0.1,0.5]$、$\sigma\in[0.5,1.5]$、$\rho_v\in[0.3,0.8]$ 全区间进入蒙特卡洛与 Tornado 扫描，结论以方向概率报告。量化结果：三条核心结论的方向概率 100\%/99.65\%/99.75\%，弹性族单参数扰动摆幅不超过 1.2\%，该误差源对结论方向的影响可接受。

\textbf{2. 构造标定误差。}需求时段份额 $\{0.18,0.42,0.16,0.24\}$、出车率与三区供求份额无官方分项统计，为构造标定。其不确定性已并入蒙特卡洛的 $D_0$ 扰动（$\pm 20\%$ 覆盖份额估计的可能偏差），Tornado 显示 $D_0$ 虽是第一敏感参数（摆幅约 32\%），但它只同比例放大或缩小成本水平，不改变时段间与方案间的相对比较，对排序结论无影响。

\textbf{3. 结构简化误差。}时段独立性假设忽略了跨时段需求转移。专项对比检验（零转移对 10\% 转移率）显示：S2、S3 场景的高峰候车估计差异分别为 9.6\% 与 5.7\%，在可接受范围；S5 场景差异 21.5\%，超过 10\% 的容忍线，原因是 S5 的高峰时段差异化补贴（高峰 11 元、其他 5 元）人为制造了跨时段套利空间。该误差只影响 S5 单方案的数值精度，不影响方案排序与方向结论，在结果解读时已按保守方向对待。

\textbf{4. 口径与边界误差。}一是 Jensen 偏差：蒙特卡洛均值与基准点值的差异（信息效应 26.2\%、资金效应 92.3\%）是 $E[f(X)]\ne f(E[X])$ 的口径差异，已通过“均值加区间为主、点值对照”的报告口径消解，非模型错误。二是封闭供给池边界：补贴期社会车辆涌入未计入，问题二的补贴效果为保守下界，上述方向判断只会被强化。三是 $\kappa$ 折算口径：平台空驶率改善按保守幅度折算，信息效应估计偏保守。这三类边界误差方向一致地使本文结论偏保守，可接受。

\subsection{鲁棒性验证}

除上述逐类误差的量化外，做两项整体鲁棒性检验。第一，极端弹性复核：需求弹性取区间强端 $\varepsilon_g=-0.9$ 重跑问题三，最优结构仍为纯司机侧纯高峰，方向一致。第二，双模型互证：排队模型与均衡模型对高峰过载、平峰宽松的判断完全一致（见问题一模型验证）。综合 Tornado 低敏、蒙特卡洛方向概率、临界点安全边际与双模型互证，本文全部核心结论在参数与结构扰动下稳健。

% ================= 七、模型评价与推广 =================
\section{模型评价与推广}

\subsection{模型优势}

第一，三问统一于同一均衡引擎：测量、评价、优化共享一套方程、参数与校验，指标口径全程一致，避免了三问各自建模导致的结论不可比。第二，反事实识别策略把“补贴是否有用”从前后对比的相关性描述升级为可分解的因果结构，信息效应与资金效应的分离使“补贴反噬”这一反直觉现象获得可检验的机理（弹性序 $\lvert\varepsilon_g\rvert>\varepsilon_s$）。第三，双口径计量纪律（撮合口径与物理口径分离）防止高估平台价值，SDMI 不变与打车难缓解的分离是单一口径模型无法给出的发现。第四，非凸优化的最优性用“多起点一致性加双层网格下界加临界点扫描”的三环证据链替代不存在的解析充要条件，全部核心结论经 2000 次蒙特卡洛仿真以方向概率报告，证据标准高于单点估计。

\subsection{模型局限}

模型的边界如实列出。其一，封闭供给池：未纳入补贴大战期间社会车辆与兼职司机的涌入，问题二的补贴效果为保守下界，若放开该假设，司机侧补贴的增产能力会更强。其二，时段独立性：忽略跨时段需求转移，对采取时段差异化补贴的方案（S5）造成最高 21.5\% 的高峰候车估计偏差，方向结论不受影响但单方案数值需谨慎引用。其三，构造标定：需求时段份额、出车率与三区份额为构造标定，虽已用蒙特卡洛传播其不确定性，绝对水平结论（如 687.2 万元）的精度仍受标定质量限制。其四，单城市静态框架：以北京日均状态为对象，未刻画周内波动、节假日与天气冲击，也未涉及网约车与巡游车的双车型竞争。

\subsection{改进方向与推广}

沿四条路径可改进。一是把供给池放开为弹性进入模型，纳入社会车辆的兼职供给，使补贴的长期效果可评估。二是引入跨时段需求转移项，把时段独立性放松为马尔可夫式转移结构。三是把单城市框架推广为多城市网络，比较不同总量管制强度下的最优补贴结构。四是把静态日均模型扩展为动态定价环境，讨论平台在高峰加价与补贴之间的政策组合。模型引擎本身可直接推广至外卖配送、即时零售等同样具有双边撮合结构的市场，只需替换服务率、时段划分与成本参数的标定。

\subsection{题目设计合理性讨论}

从建模者视角看，本题是一道设计精良的赛题。第一，无附件的开放式数据环境考察检索与标定能力，而北京案例的公开锚点丰富（总量、成交、候车、空驶率均有官方或权威报告口径），使“锚点外检”的标定验证成为可能。第二，三问递进自然：测量为评价提供基准，评价为设计识别通道，任何一问的方法选择都被前后问约束，割裂式解题难以自洽。第三，指标设计开放：题目不指定“匹配程度”的定义，迫使参赛者在撮合效率与物理配置之间做出概念辨析，双口径的必要性正是题目隐含的考察点。第四，政策含义真实：2014 年补贴大战是发生过的自然实验，滴滴、快的各轮方案参数可考，反事实评价的结论可以接受现实直觉的检验。第五，隐含约束识别空间充足：总量管制、高峰官方定义、司机参与约束、公平性约束等均需从背景材料中挖掘，考验审题深度。

% ================= 参考文献 =================
\begin{thebibliography}{99}
\zihao{-5}

\bibitem{ref1} 北京市发展和改革委员会. 关于调整本市出租汽车价格的通知: 京发改〔2013〕1052号[Z]. 北京, 2013.

\bibitem{ref2} 中华人民共和国住房和城乡建设部. 城市道路交通规划设计规范: GB 50220-95[S]. 北京: 中国建筑工业出版社, 1995.

\bibitem{ref3} 罗兰贝格管理咨询公司. 移动互联下的城市综合出行变革[R]. 上海, 2015.

\bibitem{ref4} 北京交通发展研究中心. 北京市交通发展年度报告[R]. 北京, 2014.

\bibitem{ref5} 北京市统计局, 国家统计局北京调查总队. 北京市2013年国民经济和社会发展统计公报[R]. 北京, 2014.

\bibitem{ref6} 北京市人民政府. 关于加强出租汽车管理提高运营服务水平的意见: 京政发〔2013〕13号[Z]. 北京, 2013.

\bibitem{ref7} 蒲长廷. 北京出租: 下周一打车2.3元/公里[N]. 法制晚报, 2013-06-06.

\bibitem{ref8} 亿邦动力网. 滴滴打车取消补贴: 历经这般“酸甜苦辣”[N/OL]. 2014-06-10.

\bibitem{ref9} Rochet J C, Tirole J. Platform competition in two-sided markets[J]. Journal of the European Economic Association, 2003, 1(4): 990-1029.

\bibitem{ref10} Armstrong M. Competition in two-sided markets[J]. The RAND Journal of Economics, 2006, 37(3): 668-691.

\bibitem{ref11} Douglas G W. Price regulation and optimal service standards: the taxicab industry[J]. Journal of Transport Economics and Policy, 1972, 6(2): 116-127.

\bibitem{ref12} Cairns R D, Liston-Heyes C. Competition and regulation in the taxi industry[J]. Journal of Public Economics, 1996, 59(1): 1-15.

\bibitem{ref13} Yang H, Wong S C. A network model of urban taxi services[J]. Transportation Research Part B: Methodological, 1998, 32(4): 235-246.

\bibitem{ref14} Gross D, Harris C M. Fundamentals of queueing theory[M]. 4th ed. New York: John Wiley \& Sons, 2008.

\bibitem{ref15} 梁希喆. 北京市出租车需求价格弹性分析: 基于2001\textasciitilde 2011年的年度数据[J]. 交通财会, 2013(4): 69-72.

\end{thebibliography}

% ================= 附录 =================
\appendix
\ctexset{
  section = {
    name = {附录},
    number = \Alph{section},
    format = \zihao{3}\heiti\centering,
    aftername = \quad
  }
}

\section{支撑材料文件列表}

\begin{table}[H]
\centering
\caption{支撑材料文件清单}
\zihao{-5}
\begin{tabular}{p{5.2cm}p{8.4cm}}
\toprule
文件 & 内容说明 \\
\midrule
project/equilibrium\_engine.py & 分时段双边市场广义成本均衡引擎（全部三问共享），含 M/M/c 排队论交叉验证与标定自检 \\
project/问题1\_求解.py & 问题一求解脚本：时空双维度指标计算与图 2、图 3 生成 \\
project/问题2\_求解.py & 问题二求解脚本：八场景反事实评价、蒙特卡洛 2000 次仿真与图 4、图 5 生成 \\
project/问题3\_求解.py & 问题三求解脚本：多起点 SLSQP 优化、双层网格验证、预算扫描与图 6、图 7 生成 \\
project/验证\_蒙特卡洛与排队论.py & 模型验证脚本：排队论互验、批次收敛与 Tornado 灵敏度图 8、图 9 生成 \\
project/打包\_solution.py & 结果打包脚本：汇总各问结果为 solution.json \\
result\_bundle/solution.json & 三问完整求解结果（目标值、决策变量、最优性证据、灵敏度、临界点、蒙特卡洛统计） \\
result\_bundle/figures/*.pdf & 全部插图矢量版（图 1 至图 9） \\
\bottomrule
\end{tabular}
\end{table}

\section{源程序代码}

\subsection{均衡引擎（equilibrium\_engine.py）}
\VerbatimInput{project/equilibrium_engine.py}

\subsection{问题一求解（问题1\_求解.py）}
\VerbatimInput{project/问题1_求解.py}

\subsection{问题二求解（问题2\_求解.py）}
\VerbatimInput{project/问题2_求解.py}

\subsection{问题三求解（问题3\_求解.py）}
\VerbatimInput{project/问题3_求解.py}

\subsection{模型验证（验证\_蒙特卡洛与排队论.py）}
\VerbatimInput{project/验证_蒙特卡洛与排队论.py}

\subsection{结果打包（打包\_solution.py）}
\VerbatimInput{project/打包_solution.py}

\end{document}
